% Insert for final_draft_SDE: deformation-field computation and training.
% Suggested placement:
% 1) Replace the current Section 3.3 with Sections 3.3--3.3.4 below.
% 2) Keep Sections 3.4--3.5, but replace them with the clarified versions below.
% 3) Insert Section 3.6 before the current Section 4.

\subsection{Neural SDE for Cardiac Deformation}
\label{sec:neural_sde_deformation}

Let $\Omega \subset \mathbb{R}^{3}$ be the image domain and let
$X=\{x_i\}_{i=1}^{N_v}$ denote the voxel grid. We choose a reference
cardiac frame $I_0$ (typically end-diastole) and define the deformation
at time $t_k$ as a backward transformation
\[
    \phi_k : \Omega \rightarrow \Omega,
    \qquad
    \widehat I_k(x) = I_0(\phi_k(x)).
\]
This convention is used because image warping is differentiable by
standard grid sampling: the predicted image at time $t_k$ is obtained by
sampling the reference image at the coordinates $\phi_k(x)$. We write
\[
    \phi_k(x) = x + u_k(x),
\]
where $u_k:\Omega\rightarrow\mathbb{R}^3$ is the displacement field.

Directly evolving a dense deformation vector in
$\mathbb{R}^{3N_v}$ is expensive and makes covariance propagation
intractable. We therefore model cardiac motion in a compact hidden state
$h(t)\in\mathbb{R}^{m}$ and decode the dense deformation through an
output network $c_\psi$:
\[
    u_k = c_\psi(h_k), \qquad \phi_k(x)=x+u_k(x).
\]
The hidden state follows an Itô neural SDE between observed frames,
\[
    dh(t) = f_\theta(h(t),t)\,dt
          + G_\theta(h(t),t)\,dW_t ,
    \qquad dW_t\sim\mathcal{N}(0,Q_0\,dt),
\]
where $f_\theta$ models the deterministic cardiac dynamics and
$G_\theta$ controls process uncertainty. This compact formulation is the
one used for analytical covariance propagation. It is also compatible
with NODEO-style registration because the dense deformation is still
generated by a differentiable deformation decoder and is trained by
warping losses.

\paragraph{Velocity regularization.}
To connect the model to neural ODE registration, the decoder may be
implemented either as a direct displacement decoder $u_k=c_\psi(h_k)$
or as an integrated velocity decoder. In the velocity form, a smooth
velocity field is produced from the hidden state,
\[
    v_k(x) = K * a_\psi(h_k,x,t_k),
\]
where $K$ is a Gaussian smoothing kernel and $a_\psi$ is a neural field
or convolutional decoder. The deformation is then generated by
\[
    \frac{d\phi(x,t)}{dt}=v(\phi(x,t),t), \qquad \phi(x,0)=x.
\]
The direct displacement decoder is computationally simpler, while the
velocity form more closely matches NODEO and naturally encourages
temporally coherent diffeomorphic motion. In both cases, the trainable
output is the dense field $\phi_k$ and the training objective below is
unchanged.

\subsection{CVRNN Update at Observed Frames}
\label{sec:cvgruu_update_clarified}

At each observed cine frame $I_k$, the frozen CineMA encoder supplies
the latent observation
\[
    z_k = E_{\mathrm{CineMA}}(I_k),
\]
with latent covariance $\Sigma_k$ obtained by propagating the
C-UNSURE image-noise estimate through the frozen encoder. The SDE first
propagates the previous posterior state to a prior state:
\[
    \hat h'_k, P'_k =
    \mathrm{SDESolve}(\hat h_{k-1},P_{k-1},t_{k-1},t_k).
\]
The observation is then assimilated by a covariance-propagating recurrent
cell,
\[
    \hat h_k = v_\omega(\hat h'_k,z_k),
\]
and its covariance is updated by the first-order uncertainty
transformation
\[
    P_k =
    \Delta_h v_\omega\,P'_k\,(\Delta_h v_\omega)^\top
    +
    \Delta_z v_\omega\,\Sigma_k\,(\Delta_z v_\omega)^\top ,
\]
where
\[
    \Delta_h v_\omega =
    \frac{\partial v_\omega(h,z)}{\partial h}
    \bigg|_{\hat h'_k,z_k},
    \qquad
    \Delta_z v_\omega =
    \frac{\partial v_\omega(h,z)}{\partial z}
    \bigg|_{\hat h'_k,z_k}.
\]
The first term propagates dynamic or epistemic uncertainty; the second
injects aleatoric observation uncertainty learned by C-UNSURE.

\subsection{Analytical Uncertainty Propagation}
\label{sec:analytical_uncertainty_clarified}

Between frames, the hidden-state mean and covariance are propagated by
linearizing the SDE drift:
\[
    \frac{d\hat h}{dt}=f_\theta(\hat h,t),
\]
\[
    \frac{dP}{dt}
    =
    F_h P + P F_h^\top
    +
    G_h Q_0 G_h^\top,
    \qquad
    F_h =
    \frac{\partial f_\theta(h,t)}{\partial h}
    \bigg|_{\hat h}.
\]
In the calibrated version, $Q_0$ is replaced by $\bar\rho_k Q_0$, where
$\bar\rho_k$ is the running innovation calibration factor.

The dense deformation field is decoded from the posterior hidden state:
\[
    \widehat u_k = c_\psi(\hat h_k), \qquad
    \widehat\phi_k(x)=x+\widehat u_k(x).
\]
Let
\[
    C_k =
    \frac{\partial \operatorname{vec}(c_\psi(h))}
         {\partial h}
    \bigg|_{\hat h_k}
    \in \mathbb{R}^{3N_v\times m}.
\]
The deformation covariance is
\[
    R_k = C_k P_k C_k^\top .
\]
Since storing $R_k\in\mathbb{R}^{3N_v\times 3N_v}$ is infeasible, we
use either a diagonal or voxel-wise block-diagonal approximation:
\[
    \operatorname{diag}(R_k)_i = c_{k,i}^\top P_k c_{k,i},
\]
where $c_{k,i}$ is the $i$-th row of $C_k$, or
\[
    R_k^{(x)}
    =
    C_k^{(x)}P_k(C_k^{(x)})^\top
    \in\mathbb{R}^{3\times3},
\]
where $C_k^{(x)}\in\mathbb{R}^{3\times m}$ contains the three Jacobian
rows corresponding to the displacement vector at voxel $x$. The block
form preserves local correlations between the $x$, $y$, and $z$
components of displacement while avoiding a full dense covariance.

\subsection{Training Objective for Deformation Fields}
\label{sec:deformation_training_objective}

The deformation field is learned without ground-truth displacement
labels by combining image registration losses, latent observation losses,
and deformation regularization. CineMA and C-UNSURE are frozen; the
trainable components are the SDE drift $f_\theta$, diffusion
$G_\theta$, CVRNN update $v_\omega$, observation decoder $o_\chi$, and
deformation decoder $c_\psi$.

\paragraph{Image registration loss.}
For each observed frame, the decoded deformation warps the reference
frame:
\[
    \widehat I_k = W(I_0,\widehat\phi_k),
\]
where $W$ is differentiable spatial sampling. The basic image loss is
local normalized cross-correlation,
\[
    \mathcal{L}_{\mathrm{img}}
    =
    \sum_{k=1}^{K}
    \left[
    1-\mathrm{LNCC}(\widehat I_k,I_k)
    \right].
\]
For multi-contrast or intensity-normalized cine MRI, $\mathrm{LNCC}$ can
be replaced by MSE, SSIM, or mutual information, but LNCC is the default
registration objective.

\paragraph{Latent observation likelihood.}
The hidden state also predicts the CineMA latent observation through an
observation decoder $o_\chi$:
\[
    \widehat z_k = o_\chi(\hat h'_k,t_k).
\]
With innovation $e_k=z_k-\widehat z_k$ and
\[
    S_k =
    \Delta_h o_\chi\,P'_k\,(\Delta_h o_\chi)^\top+\Sigma_k,
\]
we use the Gaussian negative log-likelihood
\[
    \mathcal{L}_{\mathrm{lat}}
    =
    \frac{1}{2}\sum_{k=1}^{K}
    \left[
    e_k^\top S_k^{-1}e_k + \log|S_k|
    \right].
\]
This term ties the SDE-RNN state to the observed CineMA trajectory and
uses the C-UNSURE covariance as the aleatoric observation noise.

\paragraph{Segmentation supervision when available.}
If ED/ES labels or sparse annotated frames are available, they are used
only as registration supervision, not as deformation ground truth. Let
$Y_0$ be the reference segmentation and $Y_k$ the target segmentation.
Then
\[
    \widehat Y_k = W(Y_0,\widehat\phi_k),
\]
and
\[
    \mathcal{L}_{\mathrm{seg}}
    =
    \sum_{k\in\mathcal{A}}
    \left[
    1-\mathrm{Dice}(\widehat Y_k,Y_k)
    \right],
\]
where $\mathcal{A}$ is the set of annotated frames. This term improves
anatomical alignment while preserving the unsupervised nature of the
deformation field.

\paragraph{Deformation regularization.}
Following NODEO-style registration, we regularize the deformation by
penalizing folding, excessive motion, and spatial roughness:
\[
    \mathcal{L}_{\mathrm{Jdet}}
    =
    \frac{1}{K N_v}
    \sum_{k,x}
    \left\|
    \operatorname{ReLU}(\epsilon-\det\nabla\widehat\phi_k(x))
    \right\|_2^2 ,
\]
\[
    \mathcal{L}_{\mathrm{mag}}
    =
    \frac{1}{K N_v}
    \sum_{k,x}
    \|\widehat u_k(x)\|_2^2 ,
\]
\[
    \mathcal{L}_{\mathrm{smt}}
    =
    \frac{1}{K N_v}
    \sum_{k,x}
    \|\nabla\widehat u_k(x)\|_2^2 .
\]
For the velocity-decoder variant, $\mathcal{L}_{\mathrm{mag}}$ is applied
to $v_k$ along the integration path. A temporal smoothness term may also
be used:
\[
    \mathcal{L}_{\mathrm{temp}}
    =
    \frac{1}{K-1}
    \sum_{k=1}^{K-1}
    \left\|
    \frac{\widehat u_{k+1}-\widehat u_k}{t_{k+1}-t_k}
    \right\|_2^2 .
\]

\paragraph{Total objective.}
The full training loss is
\[
\begin{aligned}
    \mathcal{L}_{\mathrm{total}}
    =
    &\lambda_{\mathrm{img}}\mathcal{L}_{\mathrm{img}}
    +\lambda_{\mathrm{lat}}\mathcal{L}_{\mathrm{lat}}
    +\lambda_{\mathrm{seg}}\mathcal{L}_{\mathrm{seg}} \\
    &+\lambda_{\mathrm{Jdet}}\mathcal{L}_{\mathrm{Jdet}}
    +\lambda_{\mathrm{mag}}\mathcal{L}_{\mathrm{mag}}
    +\lambda_{\mathrm{smt}}\mathcal{L}_{\mathrm{smt}}
    +\lambda_{\mathrm{temp}}\mathcal{L}_{\mathrm{temp}} .
\end{aligned}
\]
If no segmentation labels are available for a sequence, the corresponding
term is omitted by setting $\lambda_{\mathrm{seg}}=0$.

\subsection{Training Procedure}
\label{sec:training_procedure}

Training proceeds in four stages.

\begin{enumerate}
    \item \textbf{Observation preparation.}
    Each cine frame is encoded by the frozen CineMA encoder to obtain
    $z_k$. C-UNSURE estimates the image-noise kernel $\eta_k$, and the
    matrix-free finite-difference estimator produces $\Sigma_k$.

    \item \textbf{Sequence construction.}
    Frames from the same patient and view are ordered by cardiac time.
    Each training sample is
    $\{I_k,z_k,\Sigma_k,t_k\}_{k=0}^{K}$, plus optional ED/ES labels.

    \item \textbf{SDE-RNN training.}
    Starting from $h_0$, the model alternates SDE propagation and CVRNN
    updates at observed frames. At every frame it decodes
    $\widehat\phi_k$, warps the reference image, computes the image and
    latent losses, and applies deformation regularization.

    \item \textbf{Validation and calibration.}
    Validation monitors image similarity, Dice/MCD on annotated frames,
    folding rate $\det\nabla\phi_k\leq0$, and the normalized innovation
    ratio $\rho_k=d^{-1}e_k^\top S_k^{-1}e_k$. The checkpoint is selected
    by registration quality subject to stable deformation regularity and
    calibrated innovations.
\end{enumerate}

At inference, the model returns
\[
    \left\{
    \widehat\phi_k,\;
    R_k^{\mathrm{diag}}\ \mathrm{or}\ R_k^{\mathrm{block}},\;
    \widehat z_k,\;
    S_k,\;
    \bar\rho_k
    \right\}_{k=0}^{K}.
\]
The deformation mean and covariance are then propagated analytically to
clinical metrics as described in Section 4.
